\chapter*{Deutsche Zusammenfassung}
\addcontentsline{toc}{chapter}{Deutsche Zusammenfassung}

Diese Masterarbeit untersucht die Entwicklung und prototypische Umsetzung einer Telemedizin-Plattform (Shifa) für den usbekischen Versorgungskontext. Ausgangspunkt sind qualitative Befunde aus ärztlichen Interviews, die wiederkehrende Probleme in der Versorgungspraxis sichtbar machen: informelle und fragmentierte Kommunikation, manuelle Terminorganisation, uneinheitliche Dokumentation, Risiken von Datenverlust sowie Unsicherheit in rechtlich sensiblen Situationen.

Auf dieser Grundlage werden funktionale und technische Anforderungen systematisch abgeleitet und in eine implementierbare Architektur überführt. Die Lösung umfasst zwei Client-Anwendungen (Patient und Arzt/Admin), ein zentrales Backend mit rollenbasierter Zugriffskontrolle sowie persistente Datenhaltung in PostgreSQL. Ergänzend werden externe Dienste für OTP-basierte Authentifizierung, Benachrichtigungen und Videokonsultationen eingebunden.

Die Arbeit zeigt damit die Nachvollziehbarkeit des Übergangs von empirisch erhobenen Bedarfen zu konkreten Architektur- und Implementierungsentscheidungen. Gleichzeitig werden Grenzen des aktuellen Prototyps transparent ausgewiesen, insbesondere in Bezug auf produktionsreife Absicherung, Betriebsreife und vollständige Vertragsformalierung der APIs.
